% Ad Atticum 5.16 --- headnote
% ad-atticum-05-16, 11 August 51 BC, written on the road from Synnada to Philomelium

\textit{Cicero to Atticus, scribbled on the road somewhere between
Synnada and Philomelium in inland Phrygia in mid-August 51 BC ---
the Perseus dateline places the writing between the fifth and
third days before the Ides of August. Cicero has stopped on the
roadside (``I sat down on the very road'') to hand a letter to the
publicans' departing couriers and to discharge his obligation
to Atticus, who has asked for news. The opening section is a
deliberate apology for brevity; what follows is anything but
slight.}

\textit{The body of the letter is the first sustained indictment
in the correspondence of the state in which Appius Claudius
Pulcher has left the province. Cicero has now held three
successive assize-sessions --- three days at Laodicea, three at
Apamea, three at Synnada --- and the report from each is
identical: the poll-tax \textit{(epikephalaia)} cannot be paid;
the cities' rights of revenue \textit{(\=onas)} have been sold off
in advance; ``the groans, the weeping of the cities'' make of
the prior governorship not a man's work but ``a sort of
monstrousness, of some savage beast.'' Against this, Cicero
sets the discipline of his own retinue: no hay, no firewood, no
roof beyond what each man brings; for the most part they sleep
in tents, and ``the people revive at our coming.'' The closing
section flicks at Appius's flight ahead of him to Tarsus (the
furthest corner of the province), repeats the rumour --- still
just a rumour --- that Roman cavalry have been cut up by the
Parthians, and remarks dryly that Bibulus, the new governor of
neighbouring Syria, is dawdling in order to leave the later.
The march is now three days from the army camp.}
